\input{technique-note.tex}
\title{Integration Bee --- Unit 08:\\ Basic Series}
\author{UK Uni Integration Bee prep\\ Official: integration.soc.srcf.net\\ Yang Topics Guide 2nd Ed}
\date{}
\begin{document}
\maketitle
\section*{Status}
Concept teaching only. Tags: \texttt{basic-series}, \texttt{yang:1.14}.
\section{Official anchors}
\begin{itemize}
  \item Syllabus §1 --- infinite series; swapping sum and integral (convergence ignored in contest guideline).
  \item Yang §1.14 Basic Series.
\end{itemize}
\section{Learning goals}
\begin{enumerate}
  \item Integrate polynomial truncations of geometric series termwise.
  \item Connect $\int_0^{1/2}1/(1-x)\,dx$ to the geometric sum idea.
\end{enumerate}
\section{Pedagogy}
Worked $\to$ faded $\to$ interleaved. See \texttt{pedagogy/README.md}.

\section{Worked termwise --- \texttt{U08-D01} (\texttt{yang:1.14})}
\[
\int_0^1(1+x+x^2+x^3)\,dx = 1+\tfrac12+\tfrac13+\tfrac14=\tfrac{25}{12}.
\]

\section{Geometric motivation --- \texttt{U08-D02}}
\[
\int_0^{1/2}\frac{1}{1-x}\,dx=\ln 2
\]
(same value as summing $\int_0^{1/2}x^n\,dx$ over $n\ge0$).

\section{Practice}
\texttt{U08-D01}\ldots\texttt{U08-D06} in \texttt{drills.tex}; \texttt{verify/run.sh}.
\end{document}
